\documentclass[a4paper,12pt]{article}
\usepackage[T2A]{fontenc}
\usepackage[utf8]{inputenc}
\usepackage[russian]{babel}
\usepackage{amsmath,amssymb,mathtools}
\usepackage{geometry}
\geometry{left=22mm,right=22mm,top=22mm,bottom=22mm}
\usepackage{helvet}
\renewcommand{\familydefault}{\sfdefault}
\usepackage{microtype}
\usepackage{xcolor}
\usepackage{booktabs,tabularx,array}
\usepackage{hyperref}
\usepackage{tikz}
\usepackage{fancyhdr}
\usepackage{enumitem}
\usepackage{parskip}
\usepackage{setspace}
\setstretch{1.13}
\setlength{\emergencystretch}{3em}
\definecolor{accent}{HTML}{006F76}
\definecolor{darkink}{HTML}{243B41}
\definecolor{wrong}{HTML}{AD3333}
\hypersetup{colorlinks=true,linkcolor=accent,pdfauthor={Лёвин Артём Александрович},pdftitle={Ревью домашней работы — Ксения Клыкова}}
\pagestyle{fancy}
\fancyhf{}
\fancyhead[L]{\footnotesize\color{darkink}Математика · Функции по заданным точкам}
\fancyhead[R]{\footnotesize\color{darkink}26.10.2026}
\fancyfoot[C]{\footnotesize\color{darkink}Лёвин Артём Александрович · эксклюзивно для Ксении Клыковой · \thepage}
\renewcommand{\headrulewidth}{0.3pt}
\renewcommand{\footrulewidth}{0pt}
\newcommand{\err}{\textcolor{wrong}{\textbf{Ошибка:}}}
\newcommand{\idea}{\textcolor{accent}{\textbf{Ключевая идея:}}}
\newcommand{\res}{\textcolor{accent}{\textbf{Итог:}}}
\newcommand{\taskstart}[2]{\clearpage\phantomsection\label{task:#1}\section*{Задание №#1}\textcolor{accent}{\textbf{#2}}\par\medskip}
\begin{document}
\thispagestyle{empty}
\begin{center}
\vspace*{33mm}
{\Large\color{darkink}МАТЕМАТИКА}\par\vspace{18mm}
{\Huge\bfseries\color{accent}Ревью домашней работы}\par\vspace{9mm}
{\Large Линейные, квадратичные и дробно-рациональные функции}\par\vspace{18mm}
{\Large Ксения Клыкова}\par\vspace{15mm}
{\large 26 октября 2026 г.}\par\vfill
{\normalsize Лёвин Артём Александрович}\par\vspace{2mm}
{\small эксклюзивно для Ксении Клыковой}\par\vspace{17mm}
\begin{tikzpicture}[x=1cm,y=1cm]
\draw[accent!55,line width=.8pt] (-5,0) .. controls (-2.6,.55) and (-1.3,-.4) .. (0,0) .. controls (1.5,.48) and (3,-.38) .. (5,.08);
\end{tikzpicture}
\end{center}
\clearpage
\section*{Сводная таблица заданий для разбора}
\small
\renewcommand{\arraystretch}{1.45}
\begin{tabularx}{\linewidth}{@{}>{\centering\arraybackslash}p{10mm}p{26mm}X p{23mm}p{23mm}@{}}
\toprule
\textbf{№}&\textbf{Тема}&\textbf{Суть ошибки или недочёта}&\textbf{Ваш ответ}&\textbf{Правильный ответ}\\
\midrule
\hyperref[task:4]{4}&Квадратичная функция&Потерян минус при подстановке; результат совпал случайно&$4$&$4$\\
\hyperref[task:5]{5}&Квадратичная функция&Потерян знак минус перед квадратом; запись противоречива&$-5$&$-5$\\
\hyperref[task:8]{8}&Квадратичная функция&В выражении $kx^2$ пропущен квадрат аргумента&$-3$&$3$\\
\hyperref[task:9]{9}&Дробно-рациональная функция&Неправильно приведены подобные слагаемые, ответ не завершён&$b=3$&$k+a+b=9$\\
\bottomrule
\end{tabularx}
\par\normalsize
\vspace{7mm}
\noindent В таблицу включены задания, требующие исправления или завершения. Для каждого дано полное решение и одно упражнение на закрепление.
\clearpage
\thispagestyle{empty}
\vspace*{\fill}
\begin{center}
{\LARGE\bfseries\color{accent}Разбор ошибок}
\end{center}
\vspace*{\fill}

\taskstart{4}{Квадратичная функция: подстановка отрицательного аргумента}
\textbf{Текст задания.} $f(x)=x^2+bx+c$, $A(0;-2)$, $B(2;4)$. Найдите $f(-3)$.

\textbf{Ваше решение.} Правильно найдены $c=-2$ и $b=1$. При подстановке записано:
\[
f(-3)=(-3)^2+1\cdot(-3)+(-2),
\quad \text{затем}\quad 9+3-2=4.
\]
\textbf{Ваш ответ.} $4$.

\err\ Потерян минус перед тройкой при вычислении $1\cdot(-3)$. Кроме того, написанное выражение $9+3-2$ равно $10$, а не $4$.

\textbf{Где именно возникает ошибка.} Последний правильный шаг:
\[
(-3)^2+1\cdot(-3)+(-2).
\]
Первый неверный шаг --- замена произведения $1\cdot(-3)$ на $+3$. Оно равно $-3$.

\textbf{Почему это неверно.} Квадрат отрицательного числа положителен: $(-3)^2=9$. Однако второе слагаемое содержит первую степень числа $-3$; знак здесь сохраняется. Совпадение итогового ответа с правильным не подтверждает ошибочную промежуточную запись.

\idea\ Возводить в квадрат только тот множитель, над которым стоит показатель $2$, и отдельно определять знак каждого следующего слагаемого.

\textbf{Исправление от последнего правильного шага.}
\[
(-3)^2+1\cdot(-3)+(-2)=9-3-2=\boxed{4}.
\]
\textbf{Полное правильное решение.} Из $f(0)=-2$ получаем $c=-2$. Из $f(2)=4$:
\[
4=2^2+2b-2=2+2b,
\quad 2b=2,
\quad b=1.
\]
Значит, $f(x)=x^2+x-2$, поэтому $f(-3)=9-3-2=4$.

\textbf{Проверка.} $f(0)=-2$; $f(2)=4+2-2=4$.

\res\ Правильный ответ $4$, но промежуточную арифметику необходимо исправить.

\subsection*{Задача на закрепление}
Для $f(x)=x^2+bx+c$ известны $A(0;1)$ и $B(2;3)$. Найдите $f(-2)$.

\taskstart{5}{Квадратичная функция: минус перед квадратом}
\textbf{Текст задания.} $f(x)=-x^2+bx+c$, $A(0;3)$, $B(1;4)$. Найдите $f(-2)$.

\textbf{Ваше решение.} Сначала записано $4=-(1)^2+b\cdot1+c$, но в следующей строке появилось $4=1+b+c$. Далее получено $b=2$ и $f(-2)=-5$.

\textbf{Ваш ответ.} $-5$.

\err\ При переходе к следующей строке утрачен минус перед квадратом: $-(1)^2=-1$, поскольку возведение в квадрат выполняется раньше внешнего знака минус.

\textbf{Где именно возникает ошибка.} Последний правильный шаг:
\[
4=-(1)^2+b\cdot1+c.
\]
Первый неверный шаг: $4=1+b+c$; правильно $4=-1+b+c$.

\textbf{Почему это неверно.} Записи $-1^2$ и $(-1)^2$ дают разные результаты: $-1^2=-(1^2)=-1$, а $(-1)^2=1$. В функции $-x^2$ минус находится вне квадрата. Из ошибочного равенства $4=1+b+3$ следовало бы $b=0$, что противоречит следующей строке с $b=2$.

\idea\ Перед подстановкой полезно записать $-x^2=-(x\cdot x)$, чтобы сохранить внешний минус.

\textbf{Исправление от последнего правильного шага.}
\[
4=-1+b+3=b+2,\quad b=2.
\]
\textbf{Полное правильное решение.} Из $f(0)=3$ следует $c=3$. Из $f(1)=4$:
\[
4=-1+b+3,\qquad b=2.
\]
Значит,
\[
f(x)=-x^2+2x+3,
\qquad f(-2)=-(-2)^2+2\cdot(-2)+3
=-4-4+3=\boxed{-5}.
\]
\textbf{Проверка.} $f(0)=3$, $f(1)=-1+2+3=4$.

\res\ Правильный ответ $-5$; важно устранить противоречивый переход в записи.

\subsection*{Задача на закрепление}
Для $f(x)=-x^2+bx+c$ известны $A(0;2)$ и $B(1;5)$. Найдите $f(-3)$.

\taskstart{8}{Квадратичная функция: аргумент возводится в квадрат}
\textbf{Текст задания.} $f(x)=kx^2$, $A(0;0)$, $B(2;12)$. Найдите $f(-1)$ и объясните, почему одной точки $A$ недостаточно.

\textbf{Ваше решение.} Получено $12=k\cdot2^2$, затем $k=3$. В пояснении верно отмечено, что точка $A$ не позволяет найти коэффициент. Но далее записано $f(-1)=3\cdot(-1)=-3$.

\textbf{Ваш ответ.} $-3$.

\err\ При последней подстановке пропущен показатель $2$ у аргумента.

\textbf{Где именно возникает ошибка.} Последний правильный шаг --- $k=3$, поэтому $f(x)=3x^2$. Первый неправильный шаг --- вычисление по формуле $3x$ вместо $3x^2$.

\textbf{Почему это неверно.} $(-1)^2=1$: при умножении двух отрицательных единиц получается положительная единица. Следовательно, $3(-1)^2=3$, а $3(-1)=-3$. Это значения двух разных функций.

\idea\ После нахождения коэффициента сначала переписать всю функцию и только затем подставлять новый аргумент в скобках.

\textbf{Исправление и полное правильное решение.} По точке $B(2;12)$:
\[
12=k\cdot2^2=4k,\quad k=3.
\]
Отсюда
\[
f(-1)=3\cdot(-1)^2=3\cdot1=\boxed{3}.
\]
Точка $A(0;0)$ действительно не определяет $k$: для любого действительного $k$ выполняется равенство $k\cdot0^2=0$. Поэтому нужна ещё одна точка с ненулевой абсциссой.

\textbf{Проверка.} $f(2)=3\cdot4=12$, $f(0)=0$.

\res\ $f(-1)=3$. Объяснение недостаточности точки $A$ по существу верно.

\subsection*{Задача на закрепление}
Для функции $f(x)=kx^2$ известны $A(0;0)$ и $B(3;18)$. Найдите $f(-2)$ и объясните, почему точка $A$ сама по себе не определяет $k$.

\taskstart{9}{Дробно-рациональная функция: система для трёх коэффициентов}
\textbf{Текст задания.} $f(x)=\dfrac{kx+a}{x+b}$, $A(0;5)$, $B(1;4)$, $C(-2;1)$. Найдите $k+a+b$ и исключённое значение $x$.

\textbf{Ваше решение.} Верно составлены соотношения $a=5b$, $4=k+a-4b$, $-2=-2k+a-b$. После подстановки $a=5b$ записано $4=k-b$, затем $k=4+b$, $b=3$. Окончательные сумма и исключённое значение не указаны.

\err\ Неверно приведены подобные слагаемые и вследствие этого неверно найден коэффициент $b$.

\textbf{Где именно возникает ошибка.} Последний правильный шаг:
\[
4=k+5b-4b.
\]
Первый неправильный шаг:
\[
4=k-b.
\]
Правильное приведение подобных слагаемых:
\[
5b-4b=(5-4)b=b,\qquad 4=k+b.
\]
\textbf{Почему это неверно.} Коэффициент при $b$ определяется разностью $5-4=1$. Замена её на $-1$ меняет уравнение. Все дальнейшие значения становятся несовместимыми с точкой $B$.

\idea\ После умножения на знаменатели составить систему, привести подобные слагаемые, найти коэффициенты и отдельно проверить, что знаменатели в исходных точках не равны нулю.

\textbf{Исправление от последнего правильного шага.} Получаем $k=4-b$. Точка $C$ даёт $-2=-2k+4b$. Подставляем:
\[
-2=-2(4-b)+4b=-8+6b,
\quad 6b=6,
\quad b=1.
\]
Тогда $k=4-1=3$, $a=5b=5$.

\textbf{Полное правильное решение.}
Из $A(0;5)$:
\[
5=\frac{a}{b},\qquad a=5b\quad (b\ne0).
\]
Из $B(1;4)$:
\[
4=\frac{k+a}{1+b}
\ \Rightarrow\ 4+4b=k+a
\ \Rightarrow\ 4=k+a-4b.
\]
Из $C(-2;1)$:
\[
1=\frac{-2k+a}{-2+b}
\ \Rightarrow\ -2+b=-2k+a
\ \Rightarrow\ -2=-2k+a-b.
\]
Подставляем $a=5b$ в последние два уравнения:
\[
\begin{cases}
4=k+b,\\
-2=-2k+4b.
\end{cases}
\]
Отсюда $k=4-b$, $b=1$, $k=3$, $a=5$. Искомая сумма:
\[
k+a+b=3+5+1=\boxed{9}.
\]
Так как знаменатель равен $x+b=x+1$, значение $x=-1$ исключается из области определения.

\textbf{Проверка.} $f(x)=\dfrac{3x+5}{x+1}$ и $x\ne-1$.
\[
f(0)=5,\qquad f(1)=\frac82=4,
\qquad f(-2)=\frac{-6+5}{-2+1}=1.
\]
Таким образом, все три точки принадлежат графику, а их знаменатели ненулевые.

\res\ $k+a+b=9$; исключённое значение аргумента $x=-1$.

\subsection*{Задача на закрепление}
Функция $f(x)=\dfrac{kx+a}{x+b}$ проходит через точки $A(0;2)$, $B(1;3)$ и $C(-2;-6)$. Найдите $k+a+b$ и исключённое значение $x$.
\end{document}
